\subsection{Holomorphy of the algebraic key invariant in period coordinates}
\label{sec:ftb-key-period-holomorphy}

We keep the key datum and the admissible integers \(g,n,l,l'\),
the matrices \(\alpha_m^t\alpha_m=m1_4\), and the formula
\eqref{eq:ftb-kaform}.  In this subsection a complex uniformization
means the actual isomorphism of the polarized complex abelian variety
and its rigidified line bundle with its theta torus.  In particular it
identifies algebraic sections of every positive line power with the
entire functions having the multiplier specified below.  This is the
source's \(\mathrm{ThetaUnifAV}\) section comparison, together with
the specified \(\mathrm{HolUnif}\) and \(\mathrm{Wit}\); it is a
hypothesis on the \emph{same} line bundle and frame, rather than a
choice of an unrelated analytic metric.

Write
\[
 \mathfrak H_g=\{\tau=\tau^t:\operatorname{Im}\tau>0\},\quad
 [k,n]_\tau=k+\tau n,\quad
 A_\tau(k,n;z)=
 \exp\!\left(2\pi i\left(a\cdot k-b\cdot n
                      -\tfrac12 n^t\tau n-n\cdot z\right)\right).
 \tag{H1}\label{eq:ftb-level-aut}
\]
Here \(a_i,b_i\in\{0,\tfrac12\}\).
A function of level \(m\) is an entire \(F:\mathbb C^g\to\mathbb C\)
such that
\[
 F(z+[k,n]_\tau)=A_\tau(k,n;z)^mF(z).
 \tag{H2}\label{eq:ftb-level-fn}
\]
In the uniformizing frame \(\mathrm{fr}\), an algebraic section \(s\)
of the \(m\)-th line power is read as
\[
                       s(\exp z)=F(z)\,\mathrm{fr}(z)^m.
 \tag{H3}\label{eq:ftb-fnpw}
\]
The rigidification is the one used in the cubical generators, and
the ordered Lie basis corresponds to
\((2\pi i)^{-1}e_1,\ldots,(2\pi i)^{-1}e_g\).
The parameter domain in the source consists of all complex matrices
\(p\) with \(\operatorname{Im}((p+p^t)/2)>0\);
every formula below is pulled back by \(\tau=(p+p^t)/2\).
This gives the source's open domain without treating the symmetric
matrices as an open subset of all matrices.

\begin{srlemma}[Theta bases and holomorphic coordinates]
\label{lem:ftb-holo-level-coordinates}
For every \(m\geq1\) there are \(m^g\) level-\(m\) functions
\(B_{m,c}(\tau,z)\), jointly holomorphic in \((\tau,z)\), which form
a basis at every \(\tau\in\mathfrak H_g\).
If a jointly holomorphic function belongs at each parameter to the
span of a finite jointly holomorphic, pointwise independent family,
its unique coordinates are holomorphic.
The same assertions hold for external products of the theta bases.
\end{srlemma}
\begin{proof}
With \(c\) running through representatives of
\((\mathbb Z/m\mathbb Z)^g\), use the source's basis
\[
 B_{m,c}(\tau,z)=
 \theta\!\left(\tfrac cm+a,\,mb;\,mz,m\tau\right),\qquad
 \theta(u,v;z,\tau)=
 \sum_{h\in\mathbb Z^g}
 e^{\,\pi i(h+u)^t\tau(h+u)+2\pi i(h+u)\cdot(z+v)} .
 \tag{H4}\label{eq:ftb-level-theta-basis}
\]
On a compact subset of the period and \(z\) domain the imaginary
quadratic part is bounded below by a positive multiple of
\(\lVert h\rVert^2\), while the remaining part is bounded in
absolute value by a fixed multiple of \(1+\lVert h\rVert\).
A summable Gaussian therefore dominates the terms, locally
uniformly.  The same argument after multiplication by any fixed
polynomial in \(h\) handles derivatives.  Normal convergence gives
joint holomorphy.  Changing \(h\) by an integer vector gives
\eqref{eq:ftb-level-fn}.

Here is the source's dimension argument.  First take \(a=b=0\).
The integer periods give Fourier coefficients \(c_v\), \(v\in\mathbb
Z^g\); the remaining periods imply, for each coordinate vector,
\[
 c_{v+me_j}
       =\exp(\pi i m\tau_{jj}+2\pi i\,v^t\tau e_j)c_v .
\]
Thus the \(m^g\) coefficients with residues \(0\leq v_i<m\)
determine the function: if they vanish, the recurrence makes all
Fourier coefficients vanish, and Fourier uniqueness on each real
torus makes the function zero.  Consequently the dimension is at
most \(m^g\).  Translations \(z\mapsto z+k/m\) act on the nonzero
functions in \eqref{eq:ftb-level-theta-basis} by the distinct
characters \(k\mapsto e^{2\pi i c\cdot k/m}\).
Distinct character eigenvectors are independent (apply character
independence to a purported linear relation at each \(z\)).
There are \(m^g\) of them, so they form a basis.  Multiplication by
the nonzero exponential and translation
\[
 f(z)\longmapsto
 e^{\,2\pi i m(a^t\tau a/2+a\cdot(z+b))}
                         f(z+\tau a+b)
\]
is an invertible map from level \((0,0)\) to level \((a,b)\);
it sends the displayed basis to the displayed shifted basis.

For the coordinate assertion, let \(T_1,\ldots,T_N\) be the family,
and fix a parameter \(p_0\).
The evaluation vectors \((T_j(p_0,z))_j\) span \(\mathbb C^N\);
otherwise a nonzero linear functional would be a relation among
the functions.  Choose finite linear combinations of evaluations
\(\ell_i\) with \(\ell_i(T_j(p_0,-))=\delta_{ij}\).
Set \(C_{ij}(p)=\ell_i(T_j(p,-))\) and
\(b_i(p)=\ell_i(F(p,-))\).  They are holomorphic and \(C(p_0)=1\).
On \(\det C\ne0\), the coordinate vector is \(C^{-1}b\);
the adjugate formula proves holomorphy.  Uniqueness patches these
local expressions.  Independence of external-product functions
follows one variable at a time.  To obtain their spanning
assertion, expand in the first variable, use this coordinate
argument to make the coefficients holomorphic in the remaining
variables, and iterate.  Their transformation law in the
remaining variables follows from uniqueness of the first
expansion.  This is precisely the tensor-expansion argument in
\srcid{exists\_holomorphic\_coord\_simTheta}.
\end{proof}

\begin{srlemma}[The holomorphic box determinant]
\label{lem:ftb-holo-box}
Fix \(m\geq1\), a level-\(m\) basis \(B_c\), and an integral
\(\alpha\) with \(\alpha^t\alpha=m1_4\).
The algebraic unit \(\lambda_m\) of
\eqref{eq:ftb-lambda}, read in these bases, is a holomorphic
function of the period.  Its formula is the following.
Choose integer tuples \((k_I,n_I)\) representing the two lattice
coordinates of \(\ker(\alpha^t)\), and put
\[
 t_{I,r}=\frac1m\sum_c\alpha_{rc}[k_{I,c},n_{I,c}]_\tau,\qquad
 \Phi_I(y)=
 \prod_{r=1}^4\theta(a,b;(\alpha y)_r+t_{I,r},\tau)
             e^{\,2\pi i\sum_c n_{I,c}\cdot y_c}.
\]
Write, uniquely,
\[
 \Phi_I(y)=\sum_J q_{I,J}(\tau)\prod_cB_{J_c}(\tau,y_c).
\]
Then, in the source's row convention,
\[
 \lambda_m(\tau)=
 \frac{\det(q_{I,J}(\tau))^{2m^2}}
 {\displaystyle
  \left(\prod_I\prod_r
 A_\tau\!\left(\sum_c\alpha_{rc}k_{I,c},
              \sum_c\alpha_{rc}n_{I,c};0\right)\right)^2}.
 \tag{H5}\label{eq:ftb-holo-lambda}
\]
This equality has no undetermined root of unity or scalar.
\end{srlemma}
\begin{proof}
The elementary automorphy calculation using
\(\alpha^t\alpha=m1\) and the half-characteristic gives
\[
 \Phi_I(y_1,\ldots,y_c+[k,n]_\tau,\ldots,y_4)
             =A_\tau(k,n;y_c)^m\Phi_I(y).
\]
The functions are entire and jointly holomorphic in the period.
Lemma~\ref{lem:ftb-holo-level-coordinates} therefore gives the
unique holomorphic coefficients \(q_{I,J}\).

We check the algebraic normalization.  Divide an actual cubical
generator, in each line-bundle chart, by the product of its
uniformizing frames.  On overlaps both numerator and denominator
have the same transition function; their ratio is a globally
defined nowhere-zero entire function of \(y\).
For the generator \(w_I\), its product with
\(e^{2\pi i\sum n_{I,c}\cdot y_c}\) is periodic under the full
product period lattice.  Indeed the translated frame ratios
are exactly \(A_\tau\), and the exponent cancels the
similitude cross terms.  An entire function periodic under a full
real lattice is bounded on a compact fundamental parallelepiped
and hence constant.  Call this constant \(C_I\ne0\).
This argument is a quotient of actual local generators;
it does not extend a generic rational function across a boundary.

The same argument for the multiplication generator \(\mu\) shows
that its quotient by the multiplication frames is constant.
Its origin rigidification makes that constant \(1\).
If \(Q_{I,J}\) are the algebraic box coefficients of
\eqref{eq:ftb-box-generator}, comparison with the preceding
formula gives, after the permutation of kernel representatives,
\[
 q_{I,J}=C_IQ_{I,J},\qquad
 \rho_I
 \prod_r A_\tau\!\left(\sum_c\alpha_{rc}k_{I,c},
                     \sum_c\alpha_{rc}n_{I,c};0\right)
                                  =C_I^{m^2}.
\]
Insert these into
\(\lambda_m=(\det Q)^{2m^2}\prod_I\rho_I^2\).
All \(C_I\)'s cancel.  A row permutation is removed by the even
determinant exponent.  This proves \eqref{eq:ftb-holo-lambda}.
Its denominator is an exponential and never vanishes.
\end{proof}

\begin{srlemma}[Simple zero sections and varying contractions]
\label{lem:ftb-holo-zeros-torsion}
Let \(F_1(p,z),\ldots,F_g(p,z)\) be jointly holomorphic level-\(n\)
functions.  Suppose that at \(p_0\) their common zeros on the
period torus are precisely \(d\) distinct simple zeros.
After shrinking to an open neighbourhood \(W\) of \(p_0\),
there are holomorphic lifts \(z_1(p),\ldots,z_d(p)\) of all these
zeros, still distinct modulo the periods, and all simple.
If a finite graded complex \(D(p)\) has holomorphic entries,
\(D(p)^2=0\), and is contracted at \(p_0\), its squared Koszul
torsion is a nonzero holomorphic function on a possibly smaller
neighbourhood.
\end{srlemma}
\begin{proof}
The holomorphic implicit-function theorem supplies one branch at
each old zero because its \(z\)-Jacobian is invertible.
To show that they are all the zeros, use real period coordinates
\(z=x+\tau y\) with \((x,y)\in[0,1]^{2g}\).
The complement of small neighbourhoods of the old zeros is
compact.  At each of its points at least one \(F_i(p_0,-)\)
is nonzero.  A finite subcover and continuity keep at least one
nonzero for all nearby parameters.  Inside the small
neighbourhoods implicit-function uniqueness applies.
The period lattice is discrete in these real coordinates;
shrinking again keeps the branches distinct modulo periods.
Their Jacobians remain nonzero by continuity.  Periodicity
reduces every zero to this compact parallelepiped, proving the
claimed completeness of the enumeration.

For the contraction assertion fix a degree-\(+1\) matrix \(h_0\)
with \(D(p_0)h_0+h_0D(p_0)=1\).  Put
\[
 A(p)=D(p)h_0+h_0D(p).
\]
Near \(p_0\), \(\det A(p)\ne0\).  Since \(D^2=0\), \(D\) commutes
with \(A\), and hence with \(A^{-1}\).  Thus
\(h_0A^{-1}\) is a contraction.  Retain only its degree-\(+1\)
entries to obtain \(h(p)\).  This is still a contraction:
in the anticommutator only those entries can contribute to
the degree-zero part, while all other degree parts of \(1\)
are zero.  Its entries are holomorphic by the adjugate formula.
Choose once the even--odd enumeration used at \(p_0\).
The required holomorphic function is
\[
             t(p)=\det\bigl((D(p)+h(p))_{\rm even,odd}\bigr)^2.
 \tag{H6}\label{eq:ftb-holo-torsion}
\]
The determinant is nonzero because \(D+h\) is invertible;
equivalently the graded block elimination in the squared-torsion
lemma applies.  No assertion \(h(p)^2=0\) is required.
Uniqueness of squared torsion identifies
\eqref{eq:ftb-holo-torsion} with any other contraction computation.
\end{proof}

\begin{srproposition}[The same algebraic value is locally holomorphic]
\label{prop:ftb-key-holomorphic-period}
Suppose the algebraic key datum at \(p_0\) has a value
\(\operatorname{KaVal}\) for the fixed admissible
\(n,l,l',\alpha\), and is in the normalized theta uniformization
\eqref{eq:ftb-fnpw}.
There are an open neighbourhood \(W\) of \(p_0\) and a holomorphic
function \(V\) on \(W\) such that at each \(p\in W\), \emph{every}
algebraic key datum with that normalized uniformization has
a \(\operatorname{KaVal}\) equal to \(V(p)\).
Consequently the section-independent invariant, once defined at
every period by the algebraization of its polarized theta torus,
is holomorphic in period coordinates.
\end{srproposition}
\begin{proof}
Read the initial transverse section tuple as \(F_{i,0}\).
Expand it in the level-\(n\) basis at \(p_0\), and keep its
coefficients fixed as \(p\) varies.  This gives the required
jointly holomorphic \(F_i(p,z)\) of the same level.
The preceding lemma supplies all their simple zeros \(z_e(p)\).

For \(m_j=l-jn>0\) choose the level bases
\(B_{m_j,c}(p,z)\), and likewise for \(l'\).
The products \(F_iB_{m_{j+1},c}\) have level \(m_j\).
Their multiplication coordinates are holomorphic by
Lemma~\ref{lem:ftb-holo-level-coordinates}.
The evaluations \(B_{l,c}(p,z_e(p))\) are holomorphic.
Thus all entries of the augmented Koszul matrix \(D_l(p)\)
are holomorphic, and the same holds in degree \(l'\).
Their differentials square to zero because multiplication
commutes and the evaluation points are common zeros.
At \(p_0\) the algebraic exactness already proved supplies
contractions; apply \eqref{eq:ftb-holo-torsion} in both degrees.

Let \(J(p)\) be the product of the source's normalized Jacobians
of \(F_i(p,-)\) at \(z_e(p)\); its entries are the first
\(z\)-derivatives, with the fixed \(2\pi i\) Lie normalization.
Joint holomorphy implies holomorphy of these derivatives and
their evaluations, so \(J(p)\) is holomorphic and nonzero.
Lemma~\ref{lem:ftb-holo-box} supplies the \(\lambda_j(p)\) and
\(\lambda'_j(p)\).  At \(p_0\) they are nonzero by the algebraic
box determinant theorem; shrink \(W\) so that they remain so.
Use \(t_l(p),t_{l'}(p),J(p),\lambda_j(p),\lambda'_j(p)\) in
\eqref{eq:ftb-kaform}.  Since every factor is holomorphic and
nonzero, its integer powers give a holomorphic function \(V\).

It remains to verify that this analytic formula computes the
\emph{actual} algebraic value, with no surviving chart factor.
The line-power section comparison makes the theta bases
algebraic bases and makes the \(F_i\) actual algebraic sections.
If \(U=\prod_e\mathrm{fr}_{j_e}(z_e)\), the algebraic squared
torsion in degree \(l\) is the analytic squared torsion times
\(U^{2l}\): this is conjugation of its augmented evaluation
rows by \(\operatorname{diag}(\mathrm{fr}_{j_e}(z_e)^l)\).
The corresponding algebraic Jacobian is \(J(p)U^{gn}\).
The same applies in degree \(l'\).
The total exponent of \(U\) in \eqref{eq:ftb-kaform} is
\[
 8D^{3g+2}
 \{l(gn-2l')+l'(2l-gn)+gn(l'-l)\}=0.
\]
The box factors agree by \eqref{eq:ftb-holo-lambda}, and the
squared torsions agree by their uniqueness.  Therefore the
algebraic value is exactly \(V(p)\).
Section independence proved in
Proposition~\ref{prop:ftb-key-tuple-independent} removes the
particular tuple chosen to construct \(V\).
On overlapping neighbourhoods, algebraization of a theta torus
provides a test key datum at each period; both local formulae
compute its unique algebraic value, so they agree.  This patches
the holomorphic value function.

The last algebraization clause is explicit: it is the actual
polarized-theta-torus realization, denoted
\(\mathrm{TorusAlgebraic}\) in the source.  Neither this clause
nor the section comparison is deduced from holomorphy of
\(V\).  Their match with the moduli/uniformization model is
a separate geometric constructor.
\end{proof}

\paragraph{Frozen source concordance.}
The base prefix for the following paths is
\srcid{src/Development/AbelianVarietyTheory/Proofs/}.
The analytic level-space proof is
\srcid{CovolumeBounds/EqZeroOfSemicharacterMultiplier.lean}
(\srcid{d9177a8a8151f010857a1872b029ca293d2ecdd4});
jointly holomorphic bases are
\srcid{ThetaLawProofs/ExistsLevBasisReal.lean}
(\srcid{0acd91467b163727cd9ab45b449a645d19079f73}).
The fixed-evaluation coordinate proof is
\srcid{DifferentiableOnCoordOfLinearIndependent.lean}
(\srcid{bb29ce974ad27e88014fe30ce737b1ff313a9fd8}).
The tensor expansion and its holomorphic coordinates are
\srcid{ExistsHolomorphicCoordSimilitudeThetaProduct.lean}
(\srcid{222a38b04c4c65633ddc63956623b5ebc871d2ae}).
The actual frame quotient gluing is
\srcid{ExistsEntireGeneratorDivFrames.lean}
(\srcid{5115a6c11545e30287524197a677172c8334ee56});
the full-lattice periodicity calculations are
\srcid{SimilitudeGeneratorEqExpInFrames.lean}
(\srcid{767a1a7851d799760d180c9278cc4c1670a5dfe7})
and \srcid{ThetaGeneratorPullbackMulConstInFrames.lean}
(\srcid{11d7336e26d26d4ef9f784da8f389ceab34a08c2}).
The frame automorphy identity, including its extension across the
zeros of theta by the identity theorem, is
\srcid{ThetaUniformisationFrameEqAutomorphyMul.lean}
(\srcid{7e0fe6b56293d92ad64e72bacde8d9aeac55f766}).
The exact box formula is
\srcid{ExistsHolomorphicThetaOfLevelBasis.lean}
(\srcid{fa578ce841cf27bb7ad5bf5c6047b9011a4b49f2}).
Moving zeros is
\srcid{SimpleZerosMoveHolomorphically.lean}
(\srcid{f2458805d37961ce99100633147e62bb53bf388d});
varying contractions is
\srcid{ExistsDifferentiableOnTorsSq.lean}
(\srcid{f93085db0421f99f72c9d2061f17e30b7141acf2}).
Their exact assembly, including the chart-factor cancellation,
is \srcid{KoszulInvariantLocallyHolomorphicPeriodMatrix.lean}
(\srcid{34a2165e6ee4a1b7242efdac1129f8328adcddbe}).
All unqualified filenames after the first one in this paragraph
are in \srcid{ThetaLawProofs/}.
